% M25, 13.09.2026: bestehende Dokumentationsbelege in den numerischen Zitierstil überführt.
% ============================================================
% §6.3 Tools und MCP — v02 (Claude, 03.09.)
% v02 (Kollegen-Runde 1): ① „JEDER agentische Knoten erhält
%   Werkzeugsatz" -> konditional („soweit … verwenden") ② Schluss-
%   Widerspruch behoben (nicht-freigabepflichtig automatisch vs.
%   gekennzeichnet mit Freigabe; Brücke Tool Approval ≠ Human Gate)
%   ③ optional übernommen: „ausgewiesene" Werkzeuge (5.6-Gleichzug).
%   Kollegen-Urteil danach: SEMANTISCH FROZEN.
% Grundlage: kapitel-6-gesamtvertrag.md ★FROZEN (§4 Zeile 6.3,
% E43-01-defensiv) + Autor-/Kollegen-Schreibauftrag 03.09. (drei
% Leistungen: Trennung der Werkzeugausführung · Werkzeuggrenzen ·
% MCP differenzieren) + maf-doku-checks.md §Ergebnisse (V6-Merker
% 2.5: UseFunctionInvocation nur pipeline-bezogen; ChatClientAgent
% fügt die Middleware automatisch ein · MCP-SDK-Zitierweg).
% Budget ~0,75 S. KEINE Toolliste, kein Tool-Katalog.
% MCP-Alt-Spike (AgenticSdlc.McpServer/) bewusst NICHT in der Prosa
% — die Thesis hat ihn nie als Systemkomponente behauptet; er ist
% unreferenzierte Phase-1-Werkstatt (V1-Audit §4c). Negativ-Erwähnung
% würde nur Verwirrung stiften.
% ============================================================

\section{Tools und MCP}
\label{sec:realisierung-tools-mcp}

Die in Abschnitt~\ref{sec:governance-mutation} begründete Kontrolle
wird beim Werkzeugaufruf konkret: Das ausführende Programm
(\emph{Host}) legt fest, welche Funktionen einer agentischen
Komponente zur Verfügung stehen. Das Modell kann daraus einen
Aufruf wählen; die technische Ausführung übernimmt der Host. Dies setzt
die begriffliche Trennung aus den
Abschnitten~\ref{sec:tool-use-function-calling} und
\ref{sec:maf-tools-structured-output} um.

\emph{Registrierung, Auswahl und Ausführung.} Der Host registriert
die Funktionen als typisierte Werkzeuge. Für den vom Modell
gewählten Aufruf führt die Function-Invocation-Schicht der
gemeinsamen Chat-Pipeline die zugehörige Funktion aus. Diese
Ausführungsschicht ist in beide Pipeline-Varianten eingebunden;
damit steht sie auch bei Modellaufrufen ohne Agentabstraktion zur
Verfügung. Das Werkzeugergebnis oder ein Ausführungsfehler wird
als Feedback zurückgegeben.

Eine eigene Protokoll-Middleware hält Beginn, Ergebnis und Fehler
der angebundenen Werkzeugaufrufe als Ereignisse fest. So bleiben
Modellausgabe und technische Ausführung getrennt beobachtbar.
Den Bedarf zeigt der belegte Fall E43-02 aus
Abschnitt~\ref{sec:beobachtbare-ausfuehrung-informationsfuehrung}:
Der Agent behauptete, alle Pflichtartefakte erneut gelesen und
geprüft zu haben; protokolliert war nur ein erneutes Lesen eines
erzeugten Artefakts. Die Instrumentierung macht solche Abweichungen
für die erfassten Zugriffe prüfbar.

\emph{Begrenzter Werkzeugraum.} Für Stufen mit Werkzeugnutzung
begrenzt der Host den Werkzeugsatz auf die jeweilige Aufgabe. Die
Einordnungsstufe aus
Abschnitt~\ref{sec:realisierung-agent-executor-workflow} kann den
Eingang und den Bestand lesen, ihren Plan prüfen und ihn als
typisierten Operationsplan ablegen. Als kosten- oder
außenwirkungsträchtig ausgewiesene Werkzeuge gehören nicht zu
ihrem Werkzeugsatz. Wo ein Werkzeug zusätzlich menschliche
Zustimmung erfordert, wird seine Funktion um eine Freigabeprüfung
ergänzt. Den Einsatz auf der Bedienebene erläutert
Abschnitt~\ref{sec:realisierung-steward}.

\emph{Externe Anbindung über MCP.} Der Steward kann einen Teil
seiner Werkzeuge über das Model Context Protocol anbinden. Als
MCP-Client nutzt er dann den offiziellen entfernten GitHub-Server für
lesende Zugriffe \cite{github2026remotemcp}.
Die bezogenen Werkzeuge werden in derselben Werkzeugform genutzt
wie lokal registrierte Funktionen. Der Endpunktpfad beschränkt die angebotenen Werkzeuge
serverseitig auf lesende Issue-Zugriffe. MCP ergänzt damit den
kontrollierten Werkzeugraum um eine standardisierte Anbindung;
es bildet kein eigenes tragendes Architekturprinzip.

Nach der modellseitigen Auswahl werden Werkzeugaufrufe innerhalb
dieser Grenzen automatisch ausgeführt, sofern keine zusätzliche
Werkzeugfreigabe vorgesehen ist. Eine solche Zustimmung erlaubt
den einzelnen Aufruf. Über die Übernahme fachlicher Änderungen
entscheiden die Human Gates auf Workflowebene, die der folgende
Abschnitt behandelt.

% M18, 12.09.2026: Sprachliche Neufassung; historische Prüfnotizen
% aus der Vorfassung folgen als lokale Herkunftshinweise.
% [Label-Pass 03.09.: 2.4.3 = sec:maf-tools-structured-output;
%  2.2.3 = sec:tool-use-function-calling (vom Autor aus dem
%  Overleaf-Bestand eingesetzt) — TODO damit erledigt.]
% [Rückverweis-Satz (Form §7).]
% [Freeze-Pass 04.09. (Kollegen-P1-4): „Ausführungsfeedback" ist in
%  2.2.3 als RÜCKGABE (Ergebnis/Fehler) definiert — die Protokoll-
%  Middleware hält EREIGNISSE fest; Begriffe jetzt getrennt.
%  Kap.-2-Definition lokal nicht prüfbar (2.2er-Dateien nur in
%  Overleaf) — Kollegen-Zitat der eigenen Kapitel-2-Fassung.]
% [Vier-Schritte-Trennung: Modellentscheidung (FunctionCallContent)
%  · Registrierung (AIFunctionFactory.Create, z. B.
%  IngestionTools.cs:28–49) · Ausführung (UseFunctionInvocation in
%  BEIDEN Pipeline-Varianten, AgentChatPipelineBuilder.cs:44/73;
%  Microsoft.Extensions.AI-Middleware FunctionInvokingChatClient,
%  V6-Check 2.5 BESTÄTIGT) · Feedback (ToolCallLoggerMiddleware,
%  agent-seitig via AsBuilder().Use, DerivationRunner.cs:224/242;
%  Code-Kommentar: „die verlässliche Quelle für tatsächlich
%  ausgeführte Tools", ChatDecisionLoggerMiddleware.cs:51).
%  Letzter Satz = V6-Eingrenzung: ChatClientAgent fügt die
%  Middleware AUTOMATISCH ein, wenn keine vorhanden — die explizite
%  Einbindung ist deshalb Systemeigenschaft (Sicherung roher
%  IChatClient-Pfade), KEINE Voraussetzung auf Agent-Ebene; bewusst
%  ohne die falsche Behauptung „ohne sie keine Ausführung".]
% [E43-01 in der EXAKTEN Anker-Einfassung (anhang-ankertabelle:
%  „belegt die Diskrepanz zwischen Modelltext und registrierter
%  Ausführung, nicht deren Ursache oder Häufigkeit") — Gesamtvertrag
%  §4/6.3: KEINE Genese-Behauptung zur heutigen Konfiguration
%  („daraus entstand UseFunctionInvocation" wäre unbelegt).]
% [Freeze-Pass 04.09. (Kollegen-P1-5): „Speicherwerkzeuge auf …
%  Bestand" war als Core-Schreibrecht lesbar — code-wahr ist
%  save_state_change_plan = Ablegen des PLANS (Proposal before
%  Truth), kein Core-Write (IngestionTools.cs:49).]
% [v02 (Kollegen-Punkte 1 + optional): „JEDER agentische Knoten
%  erhält einen Werkzeugsatz" war unnötig universal — Werkzeugnutzung
%  ist möglicher Handlungsspielraum, kein Agentenmerkmal (Kap.-2-
%  Basis; code-wahr: der Ledger-Extraktor ist agentisch OHNE Tools,
%  SemanticLedgerExtractor = reiner IChatClient+Schema).
%  „ausgewiesene" = Terminologie-Gleichzug mit 5.6-Endform
%  („entsprechend als wirkungs- oder kostenträchtig ausgewiesene
%  Handlungen").]
% [Werkzeuggrenzen ohne Toolliste (Auftrag Punkt 2): EIN Beispiel
%  (Einordnungs-Tool-Satz, IngestionTools.cs:28–49 — nur
%  get/list/search/check/save, keine Run-/Forward-Tools).
%  Freigabe-Hülle = ApprovalRequiredAIFunction
%  (Microsoft.Extensions.AI-Typ; „does not enforce the requirement"
%  — Durchsetzung leistet der Agent/Workflow, V6-Check §10;
%  Steward-Wrapping StewardRunTools.cs:68–88 -> 6.6).
%  Merker: „Abschnitt 6.6" -> \ref GESETZT (Label-Pass 03.09.).]
% [v03 (V6-SCHLUSSPASS 03.09., Punkt 4 BESTÄTIGT): offizieller
%  Remote-Server existiert (api.githubcopilot.com/mcp/), readonly
%  per /readonly-Pfad bzw. Header, Toolsets per /x/{toolset} —
%  deckt unsere Nutzung (readonly/x/issues) und den
%  „serverseitig"-Satz.]
% [V1-Audit §4c: Steward als MCP-CLIENT produktiv
%  (StewardChatRunner.cs:31/174–183, ConnectRemoteReadonlyIssuesAsync,
%  „Governance SERVER-seitig"; ModelContextProtocol 2.1.0,
%  Host.csproj:26). „dieselbe Werkzeugform": McpClientTool erbt von
%  AIFunction — Quelle MCP-C#-SDK-Doku (V6-Check §16; Learn sagt
%  „Convert the MCP tools to AIFunction's"). Einordnungs-Schluss =
%  Auftrag Punkt 3 („standardisierte Werkzeuganbindung, kein
%  Architekturprinzip").]
% [v02 (Kollegen-Punkt 2): „Werkzeugaufrufe laufen automatisch"
%  widersprach dem eigenen Freigabe-Hüllen-Satz davor — jetzt
%  differenziert (nicht-freigabepflichtig automatisch ·
%  gekennzeichnete mit Freigabe) + saubere Brücke
%  Tool Approval ≠ Human Gate. Zwei-Ebenen-Trennung aus Matrix §C
%  (ToolApproval = Tool-Ebene, Gates = Workflow-Ebene).
%  Nummernfrei-robust.]
